\subsubsection{Kết quả quá trình huấn luyện YOLOv8n}

\paragraph{Phân phối nhãn, vị trí và kích thước đối tượng}

\figurename~\ref{fig:labels} trình bày phân bố số lượng mẫu theo từng lớp đối
tượng, đồng thời thể hiện phân bố vị trí tâm và kích thước của các hộp bao
trong tập dữ liệu. Có thể nhận thấy số lượng mẫu giữa các lớp có sự chênh lệch
đáng kể. Trong đó, lớp xe hơi chiếm tỷ lệ lớn nhất với 409,093 mẫu, tiếp theo
là lớp xe tải nhỏ với 47,412 mẫu và lớp xe buýt với 22,621 mẫu. Lớp khác chỉ
bao gồm 3\,068 mẫu, chiếm tỷ lệ rất nhỏ trong toàn bộ tập dữ liệu. Kết quả này
cho thấy tập dữ liệu tồn tại sự mất cân bằng về số lượng mẫu giữa các lớp đối
tượng.

Biểu đồ phân bố tâm hộp bao cho thấy các đối tượng xuất hiện trên nhiều vị trí
khác nhau trong ảnh, tuy nhiên mật độ phân bố không đồng đều. Phần lớn đối
tượng tập trung ở khu vực trung tâm và nửa dưới khung hình, phản ánh đặc điểm
của các cảnh giao thông khi phương tiện thường xuất hiện trên mặt đường và
trong vùng quan sát chính của camera. Mặc dù vậy, dữ liệu vẫn bao phủ tương
đối đầy đủ các vị trí trong ảnh, góp phần tăng tính đa dạng của tập huấn
luyện.

Đối với kích thước đối tượng, phân bố chiều rộng và chiều cao của hộp bao cho
thấy phần lớn các đối tượng có kích thước nhỏ. Các giá trị chiều rộng và chiều
cao tập trung chủ yếu ở vùng thấp, trong khi số lượng đối tượng có kích thước
lớn giảm dần theo xu hướng phân bố đuôi dài. Đặc điểm này có thể được giải
thích bởi vị trí lắp đặt camera thường ở khoảng cách đủ xa nhằm bao quát được
một khu vực giao thông rộng hơn. Do đó, nhiều phương tiện chỉ chiếm một phần
nhỏ diện tích ảnh, dẫn đến sự xuất hiện với tần suất cao của các hộp bao có
kích thước nhỏ trong tập dữ liệu.

Nhìn chung, tập dữ liệu có đặc trưng là phân bố nhãn không đồng đều, vị trí
xuất hiện của đối tượng tập trung chủ yếu ở khu vực trung tâm ảnh và phần lớn
đối tượng có kích thước nhỏ. Đây là những đặc điểm quan trọng cần được xem xét
trong quá trình huấn luyện và đánh giá mô hình phát hiện đối tượng.

\begin{figure}[htbp]
  \centering
  \includegraphics[width=0.75\textwidth]{labels}
  \caption{Phân bố nhãn, vị trí tâm và kích thước hộp bao của các đối tượng
  trong tập dữ liệu.}
  \label{fig:labels}
\end{figure}

\paragraph{Kết quả hội tụ của mô hình}

\figurename~\ref{results} trình bày sự thay đổi của các hàm mất mát và các chỉ
số đánh giá trong quá trình fine-tuning mô hình YOLOv8n trên tập huấn luyện và
tập xác thực. Nhìn chung giá trị các hàm mất mát trên tập huấn luyện có xu
hướng giảm dần theo số epoch mặc dù tăng nhẹ ở epoch thứ 15. Điều này cho thấy
mô hình hội tụ ổn định trên tập huấn luyện.

Tương tự, giá trị các hàm mất mát trên tập xác thực cũng có xu hướng giảm dần
theo số epoch. Các giá trị này có dao động ở các epoch đầu nhưng dần ổn định ở
nhưng epoch cuối. Từ epoch 10 đến 15, các giá trị này giảm mạnh cho thấy mô
hình đã học được đặc trưng quan trọng và có khả năng tổng quát hoá tốt hơn. Sự
khác biết giữa các đường cong hàm mất mát trên tập huấn luyện và tập xác thực
không đáng kể cho thấy mô hình chưa xuất hiện dấu hiệu quá khớp nghiêm trọng.

Đối với các chỉ số đánh giá, Precision, Recall, mAP@0.5 và mAP@0.5:0.95 đều có
xu hướng tăng trong quá trình huấn luyện. ặc dù một số epoch ghi nhận sự sụt
giảm nhẹ, các chỉ số này nhanh chóng phục hồi và đạt giá trị ổn định ở giai
đoạn cuối. Kết quả cho thấy khả năng phát hiện và định vị đối tượng của mô
hình được cải thiện đáng kể sau quá trình fine-tuning.

Nhìn chung, mô hình YOLOv8n đạt trạng thái hội tụ ổn định trên cả tập huấn
luyện và tập xác thực, đồng thời duy trì hiệu năng tốt trong các epoch cuối
của quá trình huấn luyện.

\begin{figure}[htbp]
  \centering
  \includegraphics[width=\textwidth]{results}
  \caption{Kết quả huấn luyện và xác thực của mô hình YOLOv8n trong
  quá trình fine-tuning}
  \label{results}
\end{figure}


\paragraph{Đánh giá hiệu năng phát hiện đối tượng trên tập xác thực}

\subparagraph{Đường cong Precision}

\figurename~\ref{fig:p-curve} trình bày sự thay đổi của chỉ số Precision theo
các ngưỡng confidence khác nhau. Kết quả cho thấy Precision có xu hướng tăng
khi confidence tăng. Ở các ngưỡng confidence thấp, mô hình giữ lại nhiều hộp
bao được dự đoán, bao gồm cả những dự đoán có độ tin cậy thấp, dẫn đến số
lượng dự đoán sai lớn và làm giảm Precision. Khi confidence tăng, các hộp bao
có độ tin cậy thấp dần bị loại bỏ, giúp giảm số lượng dự đoán sai và cải thiện
độ chính xác của mô hình. Đối với toàn bộ các lớp, Precision đạt giá trị cực
đại bằng 1.00 tại ngưỡng confidence khoảng 0.973. Ngoài ra Precision của lớp
xe buýt duy trì giá trị cao hơn các lớp còn lại trong khi lớp xe tải nhỏ có
Precision thấp hơn các lớp còn lại.

\begin{figure}[htbp]
  \centering
  \includegraphics[width=0.8\textwidth]{BoxP_curve}
  \caption{Đường cong Precision theo ngưỡng confidence}
  \label{fig:p-curve}
\end{figure}

\subparagraph{Đường cong Recall}

\figurename~\ref{fig:r-curve} trình bày sự thay đổi của chỉ số Recall theo các
ngưỡng confidence khác nhau. Kết quả cho thấy Recall có xu hướng giảm khi
confidence tăng. Ở các ngưỡng confidence thấp, mô hình giữ lại hầu hết các
hộp bao được dự đoán, nhờ đó số lượng đối tượng bị bỏ sót ở mức thấp nên
Recall đạt giá trị cao. Tuy nhiên, khi confidence tăng, các hộp bao có độ tin
cậy thấp dần bị loại bỏ, làm số lượng đối tượng bị bỏ sót tăng lên và khiến
Recall giảm dần. Đối với toàn bộ các lớp, Recall đạt giá trị lớn nhất khoảng
0.98 tại ngưỡng confidence bằng 0. Tương tư như Precision, lớp xe buýt vẫn có
Recall cao hơn lớp còn lại và lớp xe tải nhỏ có Recall thấp hơn các lớp còn
lại.

\begin{figure}[htbp]
  \centering
  \includegraphics[width=0.8\textwidth]{BoxR_curve}
  \caption{Đường cong Recall theo ngưỡng confidence}
  \label{fig:r-curve}
\end{figure}

\subparagraph{Đường cong Precision--Recall}

\figurename~\ref{fig:pr-curve} trình bày đường cong Precision--Recall của mô
hình đối với từng lớp đối tượng. Kết quả cho thấy các đường cong đều nằm gần
góc trên bên phải của đồ thị, phản ánh khả năng duy trì đồng thời Precision và
Recall ở mức cao trên phần lớn miền giá trị. Đối với toàn bộ các lớp, mô hình
đạt giá trị mAP@0.5 bằng 0.953, cho thấy hiệu năng phát hiện đối tượng rất tốt
trên tập xác thực.

Xét theo từng lớp, lớp xe buýt đạt giá trị AP cao nhất với 0.989, cho thấy khả
năng nhận dạng và định vị đối tượng của mô hình đối với lớp này là tốt nhất.
Tiếp theo là lớp xe hơi với AP đạt 0.970. Trong khi đó, lớp khác và xe tải nhỏ
đạt AP lần lượt là 0.936 và 0.915, thấp hơn so với các lớp còn lại, cho thấy
mô hình vẫn gặp một số khó khăn trong việc phát hiện và phân biệt các đối
tượng thuộc hai lớp này. Tuy nhiên, các giá trị AP đều lớn hơn 0.9, chứng tỏ
mô hình đạt hiệu năng phát hiện ổn định và có khả năng tổng quát hóa tốt trên
tất cả các lớp đối tượng.

\begin{figure}[htbp]
  \centering
  \includegraphics[width=0.8\textwidth]{BoxPR_curve}
  \caption{Đường cong Precision--Recall của mô hình}
  \label{fig:pr-curve}
\end{figure}

\subparagraph{Đường cong F1-score}

\figurename~\ref{fig:f1-curve} trình bày sự thay đổi của chỉ số F1-score theo
các ngưỡng confidence khác nhau. Kết quả cho thấy F1-score của mô hình tăng
nhanh khi confidence tăng từ 0 đến khoảng 0.2. Ở các ngưỡng confidence thấp,
mô hình giữ lại phần lớn các hộp bao được dự đoán, bao gồm cả những dự đoán có
độ tin cậy thấp, dẫn đến số lượng dự đoán sai tăng lên và làm giảm Precision
mặc dù Recall vẫn được duy trì ở mức cao. Khi confidence tăng từ khoảng 0.2
đến 0.8, các dự đoán có độ tin cậy thấp dần bị loại bỏ, giúp cải thiện
Precision trong khi Recall vẫn giữ ở mức tương đối cao, từ đó F1-score đạt giá
trị lớn và ổn định trong khoảng này. Khi confidence tiếp tục tăng từ khoảng
0.8 đến 1, F1-score bắt đầu giảm dần do ở các ngưỡng confidence cao, nhiều hộp
bao chứa đối tượng thực sự cũng bị loại bỏ khỏi kết quả dự đoán, làm Recall
giảm mạnh và kéo theo sự suy giảm đáng kể của F1-score.

Đối với toàn bộ các lớp, F1-score đạt giá trị cực đại khoảng 0.90 tại ngưỡng
confidence xấp xỉ 0.572, cho thấy đây là ngưỡng mang lại sự cân bằng tốt nhất
giữa Precision và Recall. Xét theo từng lớp đối tượng, tương tự, lớp xe buýt
đạt F1-score cao nhất trên phần lớn miền confidence, cho thấy mô hình có khả
năng phát hiện và phân loại lớp đối tượng này tốt hơn so với các lớp còn lại.
Ngược lại, lớp xe tải nhỏ F1-score thấp nhất, phản ánh mức độ khó khăn lớn hơn
trong quá trình nhận dạng và phân biệt lớp đối tượng này.

\begin{figure}[htbp]
  \centering
  \includegraphics[width=0.8\textwidth]{BoxF1_curve}
  \caption{Đường cong F1-score theo ngưỡng confidence}
  \label{fig:f1-curve}
\end{figure}

\subparagraph{Ma trận nhầm lẫn chuẩn hóa}

\figurename~\ref{fig:confusion-matrix} trình bày ma trận nhầm lẫn chuẩn hóa
của mô hình trên tập dữ liệu xác thực. Các giá trị trên đường chéo chính đều ở
mức cao, cho thấy phần lớn các đối tượng được phân loại đúng vào lớp tương
ứng. Cụ thể, lớp xe buýt đạt tỷ lệ nhận dạng đúng cao nhất với giá trị khoảng
0.98, tiếp theo là lớp xe hơi với 0.97. Kết quả này phù hợp với chỉ số AP và
F1-score đã phân tích trước đó, trong đó lớp xe buýt và xe hơi đều đạt các giá
trị AP và F1-score cao hơn so với các lớp còn lại. Điều này cho thấy mô hình
có khả năng học được các đặc trưng đặc trưng của hai lớp đối tượng này và phân
biệt chúng một cách hiệu quả.

Ngược lại, lớp xe tải nhỏ và khác có tỷ lệ phân loại đúng thấp hơn, cùng đạt
khoảng 0.83. Kết quả này cũng nhất quán với các chỉ số AP và F1-score đã được
trình bày trước đó, khi đây là hai lớp có hiệu năng thấp hơn so với xe hơi và
xe buýt.

Ngoài các giá trị trên đường chéo chính, một số trường hợp nhầm lẫn giữa các
lớp cũng được thể hiện thông qua các phần tử ngoài đường chéo. Đối với lớp
xe tải nhỏ, khoảng 13\% mẫu thực tế bị dự đoán thành lớp xe hơi, cho thấy mô
hình còn gặp khó khăn trong việc phân biệt hai loại phương tiện này. Điều này
có thể xuất phát từ sự tương đồng về kích thước, hình dạng.

Ngoài ra, một tỷ lệ đáng kể các vùng nền bị mô hình dự đoán thành lớp xe hơi.
Điều này cho thấy mô hình có xu hướng ưu tiên dự đoán lớp xe hơi, một phần do
sự mất cân bằng phân bố nhãn trong tập dữ liệu. Bên cạnh đó, bộ dữ liệu còn
chứa các vùng bỏ qua với các đối tượng nhỏ hoặc bị che khuất, làm gia tăng các
trường hợp dự đoán xuất hiện trên nền trong quá trình đánh giá.

\paragraph{Nhận xét tổng kết}

Kết quả đánh giá trên tập xác thực cho thấy mô hình YOLOv8n sau quá trình
fine-tuning đạt hiệu năng phát hiện đối tượng tốt với các chỉ số Precision,
Recall, F1-score và mAP@0.5 đều ở mức cao. Mô hình nhận dạng hiệu quả các lớp
xe hơi và xe buýt, đồng thời duy trì khả năng tổng quát hóa tốt trên tập dữ
liệu xác thực. Mặc dù vẫn còn một số trường hợp nhầm lẫn giữa xe hơi và xe tải
nhỏ cũng như các dự đoán sai trên vùng nền, nhìn chung mô hình đã hội tụ ổn
định và đạt được hiệu năng tương đối tốt.

\begin{figure}[htbp]
  \centering
  \includegraphics[width=0.75\textwidth]{confusion_matrix_normalized}
  \caption{Ma trận nhầm lẫn chuẩn hóa của mô hình}
  \label{fig:confusion-matrix}
\end{figure}

\subsubsection{Kết quả trên tập kiểm tra của mô hình YOLOv8n}

\tablename~\ref{tab:test_result} trình bày các chỉ số đánh giá hiệu năng của
mô hình trên tập kiểm tra. Mô hình đạt Precision bằng 0.6952 và Recall bằng
0.6171. Điều này cho thấy mô hình có khả năng nhận diện đối tượng tương đối
tốt, với 69.52\% các dự đoán là chính xác và phát hiện được 61.71\% tổng số
đối tượng thực tế trong tập kiểm tra.

Ngoài ra, mô hình đạt mAP@50 bằng 0.6241 và mAP@50--95 bằng 0.4710. Kết quả
này cho thấy mô hình có khả năng định vị và phân loại đối tượng khá hiệu quả.
Tuy nhiên, sự chênh lệch giữa hai chỉ số mAP cho thấy độ chính xác định vị của
mô hình giảm khi yêu cầu về ngưỡng IoU trở nên nghiêm ngặt hơn.

Nhìn chung, mô hình đạt hiệu năng tương đối tốt trên tập kiểm tra và cho thấy
khả năng tổng quát hóa hiệu quả đối với dữ liệu chưa xuất hiện trong quá
trình huấn luyện.

\begin{table}[h]
\centering
\caption{Kết quả đánh giá mô hình trên tập kiểm tra}
\label{tab:test_result}
  \begin{tblr}{
    colspec = {l|c},
    hline{2},
    vline{2}
  }
  Chỉ số & Giá trị \\
  Precision & 0.6952 \\
  Recall & 0.6171 \\
  mAP@50 & 0.6241 \\
  mAP@50--95 & 0.4710 \\
  \end{tblr}
\end{table}
